% New refinement. Root and path_stability independently audited all steps and constants.
\section{Conditioning on terminal data block by block}
\label{sec:second-order-block-conditioning}

A block is a maximal two-connected subgraph, with each bridge counted as a block.
An isolated vertex is a degenerate block when needed.
The cycle space of a graph is the direct sum of its block cycle spaces.
The block-cut graph of each connected component is a tree.
Its nodes are the blocks and the articulation vertices, with incidence as adjacency.

\begin{lemma}[Weighted conditional block estimate]\label{lem:weighted-block-conditioning}
Let $Q$ be a connected loopless multigraph, and let $S$ be a nonempty proper vertex set.
For each vertex $v$, let $L_v$ be its local edge-coordinate space in $\mathcal Z(Q)^*$.
Put
\[
 U=\sum_{s\in S}L_s,\qquad B=|\delta_Q(S)|.
\]
Let $\mathcal V\subseteq V(Q)\setminus S$ consist of vertices with $L_v\cap U\ne0$.
Write $I_v$ for the indicator that $v$ is closed, and put
$p_v=\mathbb P(I_v=1)=2^{-\dim L_v}$.
For arbitrary $0\le\eta_v\le1$, there is a set $E\subseteq\mathcal V$ with $|E|\le5B$ such that
\begin{equation}\label{eq:weighted-block-conditioning}
 \mathbb E\prod_{v\in\mathcal V}(1-\eta_v I_v)
 \le\frac12+\frac12\exp\!\left(-\sum_{v\in\mathcal V\setminus E}\eta_v p_v\right).
\end{equation}
\end{lemma}

\begin{proof}
Put $H=Q-S$.
In its block-cut forest, attach one distinct terminal leaf for each edge of $\delta_Q(S)$.
Attach that leaf to the articulation node of its endpoint in $H$, if that endpoint is an articulation.
Otherwise attach it to the unique block containing its endpoint.
Thus terminals at nonarticulation vertices are also retained.
Prune each tree to the union of paths between terminal leaves.
Components with fewer than two terminals contribute nothing.
Call the resulting forest $T$.
It has at most $B$ leaves.

Exclude the vertices of $\mathcal V$ adjacent to $S$, of which there are at most $B$.
The cut-space argument in Lemma~\ref{lem:local-terminal-intersection} applies to every remaining vertex.
It shows that such a vertex is an articulation separating at least two terminal-bearing components of $H-v$.
Its articulation node therefore has degree at least two in $T$.

Also exclude articulation nodes of degree at least three in $T$.
There are at most $B$ of these.
Finally exclude articulation nodes adjacent in $T$ to a block node of degree at least three.
The sum of the degrees of all branching nodes in a forest with at most $B$ leaves is at most $3B$.
This last exclusion therefore removes at most $3B$ further vertices.
Let $E$ be the union of the excluded vertices.

Each remaining vertex $v$ has exactly two terminal-bearing incident blocks in $H$.
All other incident blocks lie in branches containing no terminal.
The two terminal sides give the same nonzero functional $a_v\in U$.
The cut-space argument also shows
\begin{equation}\label{eq:block-intersection-line}
 L_v\cap U=\{0,a_v\}.
\end{equation}
Indeed, a local cut can only select a union of the components of $H-v$.
Terminal-free components contribute zero.
The two terminal-bearing components have complementary terminal sets, whose total functional is zero.

Join two remaining vertices when they belong to a common block of $H$.
Any shared block lies on the terminal forest between their articulation nodes.
It has degree two there, since neighbors of branching block nodes were excluded.
The resulting conflict graph is therefore a forest of paths.
Its two color classes contain no pair of vertices incident with the same block.
Choose a color class $J$ with
\[
 \sum_{v\in J}\eta_vp_v
 \ge\frac12\sum_{v\in\mathcal V\setminus E}\eta_vp_v.
\]

Condition on the evaluation of all functionals in $U$.
This determines every edge incident with $S$.
The remaining edge set in $H$ is a uniform member of an affine translate of $\mathcal Z(H)$.
Its random choices in different blocks are independent.

For clarity, zero-star feasibility can also be checked within each block.
At an articulation $v$, each incident block has a prescribed star parity, determined by the terminal data on its side.
The terminal-free sides have parity zero.
The other two sides have parity $a_v$.
Deleting $v$ from a nondegenerate block leaves it connected.
Hence, whenever the prescribed parity at $v$ is zero, a compatible edge set avoiding its star exists in that block.
This follows by joining its remaining odd vertices inside the connected graph obtained after deletion.
For a bridge block, parity zero simply forces its edge to be absent.
Thus the full star at $v$ can be zero exactly when $a_v=0$.

Equation~\eqref{eq:block-intersection-line} gives the exact conditional probability
\[
 \mathbb P(I_v=1\mid U)=2p_v\,\mathbf1_{\{a_v=0\}}.
\]
The indicators for $v\in J$ are conditionally independent, because their incident block sets are disjoint.
Ignoring the other factors and using $1-x\le e^{-x}$ gives
\[
 \mathbb E\!\left[\prod_{v\in\mathcal V}(1-\eta_vI_v)\,\middle|\,U\right]
 \le\exp\!\left(-2\sum_{v\in J}\eta_vp_v\mathbf1_{\{a_v=0\}}\right).
\]

Put $W=\sum_{v\in J}\eta_vp_v\mathbf1_{\{a_v=0\}}$ and $A=\sum_{v\in J}\eta_vp_v$.
Each $a_v$ is a nonzero functional, so $\mathbb EW=A/2$, while $0\le W\le A$.
If $A=0$, the claimed bound is trivial.
Otherwise the chord bound for the convex function $e^{-2x}$ gives
\[
 \mathbb E e^{-2W}\le\frac12+\frac12e^{-2A}.
\]
The choice of $J$ now proves~\eqref{eq:weighted-block-conditioning}.
No independence among the functionals $a_v$ is required.
\end{proof}

\begin{theorem}[A sharper exponential cut constant]\label{thm:sharper-exponential-cuts}
Let $Q$ be a connected loopless multigraph with $k\ge1$ marked vertices of degree at most $d\ge2$.
For a uniform even subgraph, let $Z$ count the closed marked vertices.
Then
\begin{equation}\label{eq:sharper-exponential-cuts}
 \mathbb E2^{-Z}\le\frac12+2\exp\!\left(-\frac{k}{8d2^d}\right).
\end{equation}
The same bound holds for simultaneous forest restrictions to $k$ disjoint connected cyclic regions with boundaries of size at most $d$.
\end{theorem}

\begin{proof}
A zero local space gives an expectation at most one half.
A direct sum of $2^d$ local spaces gives an expectation below one half, as before.
Otherwise choose a maximal direct-sum family with vertex set $S$ and size $t<2^d$.
Every other marked local space meets $U=\sum_{s\in S}L_s$ nontrivially.

Suppose first that $k\ge2(5d+1)2^d$.
Then $S$ is nonempty and proper.
Its boundary has size $B\le dt$.
Apply Lemma~\ref{lem:weighted-block-conditioning} with $\eta_v=1/2$ to the marked vertices outside $S$.
At least $k-t-5B$ marked vertices remain outside its exceptional set.
Each has $p_v\ge2^{1-d}$, so
\[
 \mathbb E2^{-Z}-\frac12
 \le\frac12\exp\!\left(-2^{-d}(k-t-5B)\right)
 \le\frac12\exp\!\left((5d+1)-\frac{k}{2^d}\right)
 \le\frac12\exp\!\left(-\frac{k}{2^{d+1}}\right).
\]
This is stronger than~\eqref{eq:sharper-exponential-cuts}.

If $k<2(5d+1)2^d$, then
\[
 \frac{k}{8d2^d}<\frac{5d+1}{4d}\le\frac{11}{8}<\ln4.
\]
The right side of~\eqref{eq:sharper-exponential-cuts} therefore exceeds one.
This proves the bound for all $k$.
The cyclic-region statement follows from the same quotient and internal-cycle-space argument as before.
\end{proof}

\paragraph{Weighted scope.}
Lemma~\ref{lem:weighted-block-conditioning} retains the exact closure probabilities $p_v$.
It loses only the weighted mass on its exceptional vertices and one half of the remaining mass in the path coloring.
Its bound depends on no maximum degree after the exceptional set has been chosen.
This identifies a possible route to a weighted estimate for regions with unequal boundaries.
